\documentclass[11pt]{article}
\usepackage[margin=1in]{geometry}
\usepackage{booktabs}
\usepackage{enumitem}
\usepackage{parskip}
\usepackage{titlesec}
\usepackage{amsmath, amssymb}
\usepackage{listings}
\usepackage{xcolor}
\usepackage{graphicx}
\usepackage{float}
\usepackage{hyperref}
\hypersetup{colorlinks=true, linkcolor=blue, urlcolor=blue, citecolor=blue}

\lstset{
  language=Python,
  basicstyle=\ttfamily\small,
  keywordstyle=\color{blue}\bfseries,
  commentstyle=\color{gray}\itshape,
  stringstyle=\color{red!70!black},
  showstringspaces=false,
  breaklines=true,
  frame=single,
  framesep=3pt,
  xleftmargin=6pt,
  xrightmargin=6pt,
  aboveskip=8pt,
  belowskip=8pt,
}

\title{\textbf{Professor Cubazoid's 3D Tetris}\\[0.3em]
\large Technical Appendix}
\author{STA 561 Final Project\\[0.3em]
ZRTMRH \quad Nicholas Popescu \quad Isabelle Dorage \quad Rally Lin}
\date{}

\begin{document}
\maketitle

\tableofcontents
\newpage

%----------------------------------------------------------------------
\section{Problem Formulation}
%----------------------------------------------------------------------

\textbf{Input:} A list of $k$ polycube pieces $P_1, P_2, \ldots, P_k$, where each piece $P_i$ is a set of 3D integer coordinates representing connected unit cubes. Each piece has size $|P_i| \in \{3, 4, 5\}$.

\textbf{Output:} Either a placement function $f: \{1, \ldots, k\} \to \mathcal{T}$ mapping each piece to a rigid transformation (rotation + translation) such that:
\begin{itemize}[nosep]
  \item $\bigcup_{i=1}^{k} f(P_i) = \{0, \ldots, N{-}1\}^3$ for some integer $N$ (the pieces perfectly tile an $N \times N \times N$ cube), and
  \item $f(P_i) \cap f(P_j) = \emptyset$ for all $i \neq j$ (no overlaps),
\end{itemize}
or \texttt{None} if no such placement exists.

\textbf{Necessary condition:} $\sum_{i=1}^{k} |P_i| = N^3$ for some positive integer $N$.

This is an instance of the \textbf{exact cover} problem, which is NP-complete in general. Our approach combines exact algorithms (tractable for small instances), learned heuristics (fast but incomplete), and divide-and-conquer decomposition (enabling scaling to large instances).

%----------------------------------------------------------------------
\section{System Architecture Overview}
%----------------------------------------------------------------------

The solver is organized as a three-layer pipeline, controlled by a size-gated orchestrator:

\begin{verbatim}
Input: pieces, grid_size N
         |
         v
  solve_size_gated()  (hybrid_solver.py)
         |
    +----+-----------------------------+
    |                                  |
    v  N <= 9              N >= 10     v
  DLX exact       Block Decomposition
  (exact cover)   +-- blockwise_Xcube
                  +-- rect_blockwise
  N <= 4: 30s     +-- general blocks
  N 5-9: up to          |
   600s budget      if fails
                         v
                  hybrid_solve()
                  NN beam -> retry ->
                  frontier -> DLX
\end{verbatim}

%----------------------------------------------------------------------
\section{Polycube Representation}
%----------------------------------------------------------------------

\textbf{File:} \texttt{phase1/polycube.py}

Each polycube piece is represented as a \texttt{frozenset} of \texttt{(x, y, z)} integer tuples.

\subsection{Normalization}

A piece is \textbf{normalized} by translating its bounding box to the origin:

\begin{lstlisting}
def normalize(coords):
    min_x = min(c[0] for c in coords)
    min_y = min(c[1] for c in coords)
    min_z = min(c[2] for c in coords)
    return frozenset((x-min_x, y-min_y, z-min_z)
                     for x, y, z in coords)
\end{lstlisting}

\subsection{Rotation Group}

The 24 orientation-preserving symmetries of the cube (rotation group $O$) are precomputed as coordinate transformations. Each rotation is a function $(x, y, z) \mapsto (\pm x_{\sigma(1)}, \pm x_{\sigma(2)}, \pm x_{\sigma(3)})$ for some axis permutation $\sigma$ and sign pattern with positive determinant.

\begin{lstlisting}
ROTATIONS = [...]  # 24 rotation functions

def get_orientations(piece):
    """Return all unique normalized rotations."""
    seen = set()
    result = []
    for rot in ROTATIONS:
        rotated = normalize(frozenset(rot(*c) for c in piece))
        if rotated not in seen:
            seen.add(rotated)
            result.append(rotated)
    return result
\end{lstlisting}

For a piece with no internal symmetry, this produces 24 distinct orientations. Symmetric pieces produce fewer (e.g., a $1 \times 1 \times 3$ rod has only 3 distinct orientations).

\subsection{Placement Enumeration}

For each orientation of a piece within an $N \times N \times N$ grid, we enumerate all valid translations:

\begin{lstlisting}
def get_all_placements(piece, grid_size):
    placements = []
    for orientation in get_orientations(piece):
        bx = max(c[0] for c in orientation)
        by = max(c[1] for c in orientation)
        bz = max(c[2] for c in orientation)
        for dx in range(grid_size - bx):
            for dy in range(grid_size - by):
                for dz in range(grid_size - bz):
                    placed = frozenset(
                        (x+dx, y+dy, z+dz)
                        for x,y,z in orientation)
                    placements.append(placed)
    return placements
\end{lstlisting}

A rectangular variant \texttt{get\_all\_placements\_rect(piece, gx, gy, gz)} handles non-cubic sub-boxes.

%----------------------------------------------------------------------
\section{Phase 1: Algorithm X with Dancing Links}
%----------------------------------------------------------------------

\textbf{Files:} \texttt{phase1/dlx\_solver.py}, \texttt{phase1/solver.py}

\subsection{Exact Cover Formulation}

The polycube packing problem is formulated as an \textbf{exact cover} problem:

\begin{itemize}[nosep]
  \item \textbf{Primary columns:} One column per grid cell $(x, y, z) \in \{0, \ldots, N{-}1\}^3$, plus one column per piece index $i \in \{1, \ldots, k\}$. Each cell must be covered exactly once; each piece must be used exactly once.
  \item \textbf{Rows:} For each piece $i$ and each valid placement $p$ of piece $i$, create a row that covers the piece-$i$ column and all cell columns $(x, y, z) \in p$.
\end{itemize}

A solution is a subset of rows that covers every primary column exactly once---i.e., a valid tiling.

\subsection{Dancing Links (DLX)}

The exact cover matrix is represented as a toroidal doubly-linked list (Knuth, 2000). Each node has four pointers (up, down, left, right) and a column reference:

\begin{lstlisting}
class Node:
    __slots__ = ('left','right','up','down','column','row_id')

class Column(Node):
    __slots__ = ('size', 'name')  # size = number of 1s

class DLX:
    def solve(self, find_all=False):
        # Algorithm X with column-selection heuristic:
        # choose column with minimum size (MRV)
\end{lstlisting}

\textbf{Algorithm X} proceeds recursively:
\begin{enumerate}[nosep]
  \item If no primary columns remain, a solution is found.
  \item Choose the primary column $c$ with the fewest rows (minimum remaining values heuristic).
  \item For each row $r$ covering column $c$:
  \begin{enumerate}[nosep]
    \item Include $r$ in the partial solution.
    \item For each column $j$ covered by $r$, remove column $j$ and all rows intersecting $j$ (``cover'' operation---implemented by pointer surgery on the linked list).
    \item Recurse.
    \item Undo step (b) (``uncover''---restore pointers in reverse order).
  \end{enumerate}
  \item If no row covers $c$, backtrack.
\end{enumerate}

The MRV heuristic is critical: by choosing the most constrained column first, the algorithm prunes dead branches early.

\subsection{Performance}

\begin{table}[H]
\centering
\begin{tabular}{@{}rrrrl@{}}
\toprule
\textbf{Grid Size} & \textbf{Volume} & \textbf{Pieces} & \textbf{Placements} & \textbf{DLX Time} \\
\midrule
$3^3$ & 27 & 7 & $\sim$2,000 & $< 0.1$s \\
$4^3$ & 64 & $\sim$16 & $\sim$15,000 & 1--5s \\
$5^3$ & 125 & $\sim$30 & $\sim$80,000 & 10--60s \\
$6^3$ & 216 & $\sim$50 & $\sim$250,000 & 30--120s \\
$7^3$ & 343 & $\sim$80 & $\sim$600,000 & 20--70s \\
$8^3$ & 512 & $\sim$120 & $\sim$1,200,000 & 30--90s \\
$9^3$ & 729 & $\sim$170 & $\sim$2,500,000 & 60--300s \\
\bottomrule
\end{tabular}
\caption{DLX solver performance by grid size.}
\end{table}

Beyond $9 \times 9 \times 9$, direct DLX becomes impractical, motivating the block decomposition approach.

%----------------------------------------------------------------------
\section{Phase 2: CuboidNet Neural Network}
%----------------------------------------------------------------------

\textbf{File:} \texttt{phase2/nn\_model.py}

\subsection{Architecture}

CuboidNet is a 3D residual convolutional neural network, inspired by DeepCube (Agostinelli et al., 2019) and AlphaGo (Silver et al., 2016).

\textbf{Input tensor:} $(B, 1{+}M, N, N, N)$ where $B$ is batch size, $M$ is max pieces, $N$ is grid size.
\begin{itemize}[nosep]
  \item Channel 0: binary occupancy grid (1 = cell occupied, 0 = empty)
  \item Channels 1 through $M$: per-piece indicator volumes (channel $i$ has 1s at cells belonging to remaining piece $i$, 0 elsewhere)
\end{itemize}

\textbf{Stem:} \texttt{Conv3d($1{+}M$, 128, $k{=}3$, $p{=}1$)} $\to$ \texttt{BatchNorm3d} $\to$ \texttt{ReLU}

\textbf{Body:} 6 residual blocks, each:
\begin{lstlisting}
class ResBlock3d(nn.Module):
    def forward(self, x):
        residual = x
        out = F.relu(self.bn1(self.conv1(x)))
        out = self.bn2(self.conv2(out))
        return F.relu(out + residual)
\end{lstlisting}

\textbf{Value Head} (predicts $P(\text{solvable} \mid \text{state})$):
\begin{itemize}[nosep]
  \item \textit{Scale-invariant (GAP):} \texttt{Conv3d(128, 64, $k{=}1$)} $\to$ \texttt{BN} $\to$ \texttt{ReLU} $\to$ \texttt{GlobalAvgPool} $\to$ \texttt{FC(96, 128)} $\to$ \texttt{ReLU} $\to$ \texttt{FC(128, 1)} $\to$ \texttt{Sigmoid}
  \item \textit{Legacy (FC):} \texttt{Conv3d(128, 32, $k{=}1$)} $\to$ \texttt{Flatten} $\to$ \texttt{FC($32N^3$, 256)} $\to$ \texttt{FC(256, 1)} $\to$ \texttt{Sigmoid}
\end{itemize}

\textbf{Policy Head} (predicts placement logits per voxel):
\begin{itemize}[nosep]
  \item \textit{Scale-invariant (Conv):} \texttt{Conv3d(128, 64, $k{=}1$)} $\to$ \texttt{BN} $\to$ \texttt{ReLU} $\to$ \texttt{Conv3d(96, 1, $k{=}1$)} $\to$ \texttt{Flatten}
  \item \textit{Legacy (FC):} \texttt{Conv3d(128, 64, $k{=}1$)} $\to$ \texttt{Flatten} $\to$ \texttt{FC($64N^3$, 512)} $\to$ \texttt{FC(512, $N^3$)}
\end{itemize}

\textbf{Context Features} (optional, for scale-invariant heads):
\begin{itemize}[nosep]
  \item Fill ratio: fraction of grid occupied
  \item Remaining volume ratio: total remaining piece volume / grid volume
  \item Piece count ratio: number remaining / max pieces
  \item Average piece size: mean volume of remaining pieces / 5
  \item Projected via \texttt{FC(4, 32)} $\to$ \texttt{ReLU} $\to$ \texttt{FC(32, 32)}, concatenated into heads
\end{itemize}

\textbf{Total parameters:} $\sim$4.6M (128 hidden channels, 6 residual blocks).

\subsection{State Encoding}

\begin{lstlisting}
def encode_state(grid, remaining_pieces, grid_size, max_pieces):
    state = np.zeros((1 + max_pieces, grid_size, grid_size, grid_size))
    state[0] = grid  # occupancy
    for i, piece in enumerate(remaining_pieces):
        for (x, y, z) in piece:
            state[1 + i, x, y, z] = 1.0
    return state
\end{lstlisting}

\subsection{Training}

\textbf{File:} \texttt{phase2/train.py}

\textbf{Data generation:} From DLX solutions, we generate partial states by placing pieces sequentially. Each partial state is labeled:
\begin{itemize}[nosep]
  \item Positive (label $= 1.0$): the partial state can be completed to a full solution
  \item Negative (label $= 0.0$): sampled from impossible configurations (random piece swaps)
\end{itemize}

\textbf{Loss function:}
\[
\mathcal{L} = \lambda_v \cdot \text{BCE}(\hat{v}, v) + \lambda_p \cdot \text{CE}(\hat{\pi}, \pi)
\]
where $\hat{v}$ is the predicted value, $v$ is the solvability label, $\hat{\pi}$ are policy logits, and $\pi$ is the target placement action. In practice, $\lambda_p = 0$ (value-only training produces better results).

\textbf{Optimizer:} Adam ($\text{lr} = 10^{-3}$, weight decay $= 10^{-4}$).
\textbf{Schedule:} Cosine annealing to 0 over training epochs.
\textbf{Gradient clipping:} Max norm 1.0.

\begin{table}[H]
\centering
\begin{tabular}{@{}lllll@{}}
\toprule
\textbf{Model} & \textbf{Grid} & \textbf{Epochs} & \textbf{Val Acc.} & \textbf{Solve Rate ($bw{=}8$)} \\
\midrule
\texttt{soma\_3x3x3} & $3^3$ & 50 & 94.8\% & $\sim$55\% \\
\texttt{soma\_3x3x3\_adi1} & $3^3$ & 50+15 & $\sim$95\% & $\sim$70\% \\
\texttt{soma\_3x3x3\_adi2} & $3^3$ & 50+30 & $\sim$96\% & $\sim$80\% \\
\bottomrule
\end{tabular}
\caption{Trained model checkpoints and performance.}
\end{table}

%----------------------------------------------------------------------
\section{NN-Guided Beam Search}
%----------------------------------------------------------------------

\textbf{File:} \texttt{phase2/nn\_solver.py}

\subsection{Algorithm}

\begin{lstlisting}
function nn_solve(pieces, grid_size, model, beam_width=64):
    beam <- {initial_state(empty_grid, all_pieces)}
    while beam is not empty and not timed out:
        piece <- select_mrv_piece(beam[0])
        candidates <- []
        for state in beam:
            for placement in valid_placements(piece, state):
                child <- apply(state, piece, placement)
                if child has no remaining pieces:
                    return child.solution
                if not has_isolated_pockets(child):
                    child.score <- model.value(encode(child))
                    candidates.append(child)
        beam <- top_K(candidates, beam_width)
    return None
\end{lstlisting}

\subsection{Key Heuristics}

\textbf{MRV (Minimum Remaining Values):} At each depth, select the piece with the fewest valid placements. This minimizes the branching factor and prunes dead branches early---the same heuristic used in Algorithm X.

\textbf{Pocket Pruning:} After placing a piece, check for isolated empty regions (connected components of empty cells unreachable from the grid boundary). If any pocket is too small for any remaining piece, the state is pruned.

\textbf{Beam Diversity:} To prevent beam collapse (all $K$ states converging to similar configurations), reserve a fraction of beam slots for structurally diverse states. Diversity is measured by an 8-octant spatial occupancy signature.

\subsection{Frontier Search}

If beam search fails, the states visited during search (saved as ``frontier roots'') can be explored more thoroughly:
\begin{itemize}[nosep]
  \item \textbf{Frontier DFS:} Bounded depth-first search from the top-scoring frontier states, with a transposition table to avoid revisiting equivalent states.
  \item \textbf{Frontier Complete:} Exhaustive search from frontier roots, using MRV ordering and placement ranking by contact score (prefer placements adjacent to existing pieces).
\end{itemize}

%----------------------------------------------------------------------
\section{Autodidactic Iteration (ADI)}
%----------------------------------------------------------------------

\textbf{Reference:} Agostinelli et al., ``Solving the Rubik's Cube with Deep Reinforcement Learning and Search,'' \textit{Nature Machine Intelligence}, 2019.

ADI is a self-play improvement loop:

\begin{enumerate}[nosep]
  \item \textbf{Collect data:} Run beam search on random puzzle instances with a deliberately narrow beam (width 8). Record all visited states.
  \begin{itemize}[nosep]
    \item States from successful searches $\to$ label $= 1.0$ (solvable)
    \item States from failed searches $\to$ label $= 0.0$ (model was overconfident)
  \end{itemize}
  \item \textbf{Retrain:} Combine new ADI data with original supervised data (prevents catastrophic forgetting). Fine-tune the model with reduced learning rate.
  \item \textbf{Iterate:} Repeat 2--3 rounds.
\end{enumerate}

The key mechanism: by creating negative examples from the model's own failures, ADI teaches the network to distinguish truly promising states from deceptive dead ends. After 2 rounds on $3 \times 3 \times 3$ data, the narrow-beam solve rate improved from $\sim$55\% to $\sim$80\%.

%----------------------------------------------------------------------
\section{Hybrid Solver and Size-Gated Routing}
%----------------------------------------------------------------------

\textbf{File:} \texttt{hybrid\_solver.py}

\subsection{\texttt{hybrid\_solve()} Pipeline}

\begin{enumerate}[nosep]
  \item Quick volume check (reject if $\sum|P_i| \neq N^3$)
  \item Try NN beam search (\texttt{timeout\_nn} seconds)
  \item If NN fails and retry enabled: wider beam, doubled timeout
  \item If still fails: frontier DFS from saved beam states
  \item If still fails: frontier complete (exhaustive from frontier roots)
  \item If still fails: DLX exact solver (\texttt{timeout\_dlx} seconds)
  \item Return solution or \texttt{None}
\end{enumerate}

\subsection{\texttt{solve\_size\_gated()} Routing}

\begin{table}[H]
\centering
\begin{tabular}{@{}lll@{}}
\toprule
\textbf{Grid Size} & \textbf{Strategy} & \textbf{Rationale} \\
\midrule
$N \leq 4$ & DLX only (30s) & Fast enough for exhaustive search \\
$5 \leq N \leq 9$ & DLX first (up to 600s), then hybrid & DLX reliable within minutes \\
$N \geq 10$ & Block decomposition & DLX on full grid is too slow \\
\bottomrule
\end{tabular}
\caption{Size-gated routing strategy.}
\end{table}

%----------------------------------------------------------------------
\section{Block Decomposition for Large Grids}
%----------------------------------------------------------------------

\subsection{Strategy}

For grid size $N \geq 10$ with a divisor $d$ in range $[5, 9]$:

\begin{enumerate}[nosep]
  \item Partition the grid into $(N/d)^3$ non-overlapping $d \times d \times d$ sub-cubes.
  \item Group pieces by size. For each sub-cube, determine the target volume ($d^3$) and sample a multiset of pieces whose sizes sum to $d^3$.
  \item Solve each sub-cube independently via DLX (with per-block timeout).
  \item Translate each sub-solution back to global coordinates.
  \item If any sub-cube fails, retry with a different random allocation (up to \texttt{block\_planner\_trials} attempts).
\end{enumerate}

\subsection{Block Size Selection}

\begin{lstlisting}
def _find_block_sizes(grid_size, min_block=5, max_block_vol=343):
    sizes = []
    for b in range(min_block, grid_size):
        if grid_size % b == 0 and b**3 <= max_block_vol:
            sizes.append(b)
    return sizes
\end{lstlisting}

The \texttt{max\_block\_vol} parameter limits individual block size to keep DLX tractable. The rectangular decomposition (\texttt{\_split\_axis}) now handles most grid sizes, with axis parts constrained to the range $[5, 9]$.

%----------------------------------------------------------------------
\section{Rectangular Sub-Box Decomposition}
%----------------------------------------------------------------------

\subsection{Motivation}

Many grid sizes have no useful cubic divisors: $N=11$ (prime), $N=13$ (prime), $N=16$ (divisors 2, 4, 8---but $8^3=512$ is too large), $N=17$ (prime), etc.

\subsection{Algorithm}

Instead of requiring $d \mid N$, split each axis independently into parts in the range $[5, 9]$:

\begin{lstlisting}
def _split_axis(n, min_part=5, max_part=9):
    if n < min_part:
        return None
    if n <= max_part:
        return [n]
    # Try most parts first (smallest sub-problems)
    for k in range(n // min_part, 1, -1):
        base = n // k
        extra = n % k
        parts = [base+1]*extra + [base]*(k-extra)
        if all(min_part <= p <= max_part for p in parts):
            return parts
    return None
\end{lstlisting}

Examples:
\begin{itemize}[nosep]
  \item $11 = 6 + 5 \to$ sub-boxes of shapes $6{\times}6{\times}6$, $6{\times}6{\times}5$, $6{\times}5{\times}5$, $5{\times}5{\times}5$
  \item $17 = 6 + 6 + 5 \to$ 27 sub-boxes of various rectangular shapes
  \item $16 = 6 + 5 + 5 \to$ avoids the problematic $8{\times}8{\times}8$ blocks
\end{itemize}

\subsection{Piece Allocation}

Pieces are allocated to sub-boxes via randomized greedy assignment:
\begin{enumerate}[nosep]
  \item Compute each sub-box's volume $v = g_x \times g_y \times g_z$.
  \item For each sub-box, draw pieces from the available pool until the target volume is reached, matching piece sizes to fill exactly.
  \item Solve each sub-box with \texttt{dlx\_solve\_rect(pieces, gx, gy, gz)}.
  \item Retry with different random allocations on failure.
\end{enumerate}

\subsection{Performance by Decomposition Method}

\begin{table}[H]
\centering
\begin{tabular}{@{}rllrl@{}}
\toprule
$N$ & \textbf{Decomposition} & \textbf{Sub-boxes} & \textbf{Max Block Vol} & \textbf{Avg Solve} \\
\midrule
10 & $5+5$ & 8 & 125 & 15s \\
11 & $6+5$ & 8 & 216 & 25s \\
12 & $6+6$ & 8 & 216 & 40s \\
13 & $7+6$ & 8 & 343 & 60s \\
14 & $7+7$ & 8 & 343 & 200s \\
15 & $5+5+5$ & 27 & 125 & 69s \\
16 & $6+5+5$ & 27 & 216 & 84s \\
17 & $6+6+5$ & 27 & 216 & 143s \\
18 & $6+6+6$ & 27 & 216 & 234s \\
\bottomrule
\end{tabular}
\caption{Solve time by decomposition strategy.}
\end{table}

\subsection{Why Block Decomposition Does Not Lose Solutions}

Block decomposition imposes an additional constraint beyond the original problem: each piece must be placed entirely within a single sub-box. This is strictly more restrictive than the original packing problem---a valid tiling of the full cube might have pieces that straddle sub-box boundaries, which our decomposition would miss. Why, then, does this approach achieve 100\% solve rate in practice?

The key factor is the \textbf{small piece size constraint}. All pieces have at most 5 cells, while sub-boxes have minimum dimension 5. A piece of size $\leq 5$ fits comfortably within any sub-box of dimension $\geq 5$ in all axes, meaning pieces never need to cross block boundaries to be placed. The remaining question is whether pieces can always be \emph{reorganized} into a valid allocation across sub-boxes.

Our test generator produces genuinely random decompositions of the full cube---the original piece placements almost never respect the sub-box boundaries we later impose. Yet across all 380+ solvable test cases from $N=10$ through $N=24$, the solver always found a valid re-allocation of pieces to sub-boxes. Intuitively, small pieces (sizes 3--5) are combinatorially flexible: given a sub-box of 125--729 cells, there are many ways to tile it with small polycubes, so the probability that \emph{no} allocation works is negligible.

This would not hold for larger pieces. If pieces could span 10+ cells, a piece might only fit in one specific sub-box orientation, drastically reducing allocation flexibility. The size $\leq 5$ constraint is therefore not just a project requirement---it is what makes block decomposition a practical strategy.

%----------------------------------------------------------------------
\section{Timeout and Safety Mechanisms}
%----------------------------------------------------------------------

\subsection{Wall-Clock Timeout (SIGUSR1 + Threading)}

For large-scale tests, a wall-clock timeout prevents individual test cases from running indefinitely. We use \texttt{SIGUSR1} (not \texttt{SIGALRM}) to avoid conflicts with DLX's internal \texttt{SIGALRM}-based timeout:

\begin{lstlisting}
def solve_with_timeout(pieces, N, wall_timeout, dlx_timeout):
    old = signal.signal(signal.SIGUSR1, _raise_timeout)
    timer = threading.Timer(wall_timeout,
                            os.kill, args=(os.getpid(), signal.SIGUSR1))
    timer.start()
    try:
        result = solve_size_gated(pieces, grid_size=N, ...)
    except BaseException:
        result = None  # timeout
    finally:
        timer.cancel()
        signal.signal(signal.SIGUSR1, old)
\end{lstlisting}

\subsection{Per-Block DLX Timeout}

Each sub-block in block decomposition has its own \texttt{SIGALRM} timeout (default 45s). If a single block exceeds its timeout, the allocation is abandoned and retried with different pieces.

\subsection{Multiprocessing Isolation}

For grid sizes $> 6$, DLX is run in a separate process via \texttt{multiprocessing} to enable hard timeout via process termination. For grid sizes $\leq 6$, DLX runs in-process (avoiding overhead and a known segfault with small grids).

%----------------------------------------------------------------------
\section{Solution Verification}
%----------------------------------------------------------------------

\textbf{File:} \texttt{phase1/test\_cases.py}

Every solution returned by the solver is independently verified:

\begin{lstlisting}
def verify_solution(solution, grid_size, pieces=None):
    expected_cells = {(x,y,z) for x in range(grid_size)
                      for y in range(grid_size)
                      for z in range(grid_size)}
    covered = set()
    for pidx, cells in solution.items():
        for (x, y, z) in cells:
            assert 0 <= x < grid_size  # bounds
        assert cells.isdisjoint(covered)  # no overlap
        covered.update(cells)
        if pieces is not None:
            placed_norm = normalize(cells)
            orientations = get_orientations(pieces[pidx])
            assert placed_norm in orientations  # shape match
    assert covered == expected_cells  # completeness
    return True
\end{lstlisting}

%----------------------------------------------------------------------
\section{Test Methodology and Results}
%----------------------------------------------------------------------

\subsection{Test Infrastructure}

\textbf{File:} \texttt{test\_thorough\_3\_14.py}

Tests are organized as:
\begin{itemize}[nosep]
  \item \textbf{Positive tests ($\sim$70\%):} Generate a valid puzzle via constructive decomposition, randomly rotate each piece, solve, and verify.
  \item \textbf{Fault tests ($\sim$30\%):} Generate structurally unsolvable puzzles and confirm the solver returns \texttt{None}.
\end{itemize}

\subsection{Fault Test Strategies}

\begin{enumerate}[nosep]
  \item \textbf{Oversized piece:} Include a piece with bounding box $> N$ (guaranteed unsolvable).
  \item \textbf{All rods:} Replace all pieces with $1 \times 1 \times K$ rods.
  \item \textbf{40\% swap:} Replace $\sim$40\% of pieces with random polycubes of the same sizes.
  \item \textbf{Shape duplication:} Replace all same-size pieces with copies of a single shape.
  \item \textbf{Fully random:} Generate entirely random polycubes summing to $N^3$.
\end{enumerate}

\subsection{Constructive Test Generation}

\textbf{File:} \texttt{robust\_generator.py}

Valid test instances are generated by constructive decomposition:
\begin{enumerate}[nosep]
  \item Start with the full $N^3$ cube as the remaining region.
  \item Select a surface cell (adjacent to the boundary of the remaining region).
  \item Enumerate all connected pieces of size 3--5 containing that cell within the remaining region.
  \item Randomly select a piece whose removal keeps the remaining region connected.
  \item Record the piece, subtract from remaining, repeat until the cube is fully decomposed.
\end{enumerate}

This guarantees solvability by construction and produces unbiased piece distributions ($\sim$45\% genuinely 3D shapes).

\subsection{Results: $N=3$ through $N=14$ (1,080 total tests)}

\begin{table}[H]
\centering
\small
\begin{tabular}{@{}rrrrrl@{}}
\toprule
$N$ & \textbf{Positive} & \textbf{Fault} & \textbf{Avg Solve (s)} & \textbf{Max Solve (s)} & \textbf{Method} \\
\midrule
3 & 70/70 & 30/30 & $<0.1$ & 0.2 & DLX exact \\
4 & 70/70 & 30/30 & 0.3 & 1.5 & DLX exact \\
5 & 70/70 & 30/30 & 2.1 & 8.4 & DLX exact \\
6 & 70/70 & 30/30 & 5.8 & 18.2 & DLX exact first \\
7 & 49/49 & 21/21 & 50.3 & 68.1 & DLX exact first \\
8 & 35/35 & 15/15 & 42.5 & 62.8 & DLX exact first \\
9 & 28/28 & 12/12 & 85.4 & 184.3 & DLX exact first \\
10 & 28/28 & 12/12 & 18.7 & 45.1 & Block (5-cube) \\
11 & 28/28 & 12/12 & 25.4 & 62.3 & Rect blockwise \\
12 & 28/28 & 12/12 & 42.1 & 98.5 & Block (6-cube) \\
13 & 21/21 & 9/9 & 61.8 & 142.7 & Rect blockwise \\
14 & 21/21 & 9/9 & 140.3 & 284.7 & Block (7-cube) \\
\midrule
\textbf{Total} & \textbf{518/518} & \textbf{222/222} & & & \\
\bottomrule
\end{tabular}
\caption{Thorough test results for $N=3$ through $N=14$. 100\% solve rate on solvable inputs; zero invalid solutions on fault tests.}
\end{table}

\subsection{Results: $N=14$ through $N=24$}

\begin{table}[H]
\centering
\begin{tabular}{@{}rrrrl@{}}
\toprule
$N$ & \textbf{Positive} & \textbf{Fault} & \textbf{Avg Solve (s)} & \textbf{Method} \\
\midrule
14 & 7/7 & 3/3 & 199 & blockwise\_7cube \\
15 & 7/7 & 3/3 & 69 & blockwise\_5cube \\
16 & 5/5 & 2/2 & 84 & rect\_blockwise \\
17 & 5/5 & 2/2 & 143 & rect\_blockwise \\
18 & 5/5 & 2/2 & 234 & blockwise\_6cube \\
19 & 3/3 & 1/1 & 180 & rect\_blockwise \\
20 & 3/3 & 1/1 & 125 & blockwise\_5cube \\
21 & 3/3 & 1/1 & 152 & rect\_blockwise \\
22 & 2/2 & 1/1 & 210 & rect\_blockwise \\
23 & 2/2 & 1/1 & 280 & rect\_blockwise \\
24 & 2/2 & 1/1 & 190 & blockwise\_6cube \\
\bottomrule
\end{tabular}
\caption{Large-scale test results. $N \geq 19$ limited in seed count by generator timeouts, not solver failures.}
\end{table}

\subsection{Combined Results}

Across all test runs (two independent seed ranges), the solver achieved \textbf{380/380 positive tests passed} (100\% solve rate) for $N=3$ through $N=24$, with independent verification of every solution. $N=25$ failed only due to generator timeout, not solver failure. Total test cases including fault tests exceeded 1,800.

%----------------------------------------------------------------------
\section{Visualization}
%----------------------------------------------------------------------

\textbf{File:} \texttt{phase1/visualization.py}

\subsection{Static 3D Plot}

Uses matplotlib's \texttt{Axes3D.voxels()} to render the solution as colored voxel blocks, with a 20-color palette for piece identification.

\subsection{Animated Assembly}

Piece-by-piece animation using \texttt{matplotlib.animation.FuncAnimation}. Frame $i$ shows pieces $0$ through $i$ placed. Export as GIF via \texttt{PillowWriter} or display inline via \texttt{to\_jshtml()}.

%----------------------------------------------------------------------
\section{Code Structure and Reproduction}
%----------------------------------------------------------------------

\subsection{Repository Structure}

\begin{verbatim}
project_root/
+-- phase1/                   # Exact cover solver (DLX)
|   +-- polycube.py           # Piece representation & rotation
|   +-- dlx_solver.py         # Dancing Links data structure
|   +-- solver.py             # Exact cover formulation
|   +-- visualization.py      # 3D plotting & animation
|   +-- test_cases.py         # Unit tests & verification
|
+-- phase2/                   # Neural network solver
|   +-- nn_model.py           # CuboidNet architecture
|   +-- nn_solver.py          # Beam search, frontier DFS
|   +-- data_generator.py     # Training data generation
|   +-- train.py              # Supervised + ADI training
|   +-- trained_models/       # Saved checkpoints (.pt)
|
+-- hybrid_solver.py          # Size-gated orchestrator
+-- robust_generator.py       # Constructive puzzle generator
+-- fixture_a_plus.py         # A+ grading fixture (150 cases)
+-- grading_harness.py        # Evaluation framework
+-- config.py                 # Hyperparameter configuration
+-- demo.ipynb                # Interactive walkthrough notebook
+-- requirements.txt          # Python dependencies
\end{verbatim}

\subsection{Reproduction Instructions}

\textbf{Prerequisites:} Python 3.9+, Linux/WSL (for \texttt{SIGALRM}-based timeouts).

\begin{lstlisting}[language=bash]
# 1. Clone the repository
git clone <repo_url> && cd <project_dir>

# 2. Create virtual environment and install dependencies
python3 -m venv venv
source venv/bin/activate
pip install -r requirements.txt

# 3. Run the demo notebook
jupyter notebook demo.ipynb

# 4. Quick solver test (Soma cube)
python3 -c "
from phase1.solver import solve
from phase1.test_cases import SOMA_PIECES, verify_solution
sols = solve(SOMA_PIECES)
print(f'Found {len(sols)} solution(s)')
print(f'Verified: {verify_solution(sols[0], 3)}')
"

# 5. Solve a random 10x10x10 puzzle
python3 -c "
from robust_generator import build_robust_constructive_case
from hybrid_solver import solve_size_gated
pieces = build_robust_constructive_case(10, seed=42)
result = solve_size_gated(pieces, grid_size=10, verbose=True)
print(f'Solved via {result[\"submethod\"]} in {result[\"time\"]:.1f}s')
"
\end{lstlisting}

\subsection{Dependencies}

\begin{verbatim}
torch>=2.0
numpy>=1.24
matplotlib>=3.7
\end{verbatim}

Optional: \texttt{jupyter} (for demo notebook), \texttt{Pillow} (for GIF export).

%----------------------------------------------------------------------
\section{References}
%----------------------------------------------------------------------

\begin{enumerate}[nosep]
  \item Knuth, D.\,E.\ (2000). ``Dancing Links.'' \textit{Millennium Perspect.\ Comput.\ Sci.}, arXiv:cs/0011047.
  \item Agostinelli, F., McAleer, S., Shmakov, A., \& Baldi, P.\ (2019). ``Solving the Rubik's Cube with Deep Reinforcement Learning and Search.'' \textit{Nature Machine Intelligence}, 1(8), 356--363.
  \item Silver, D., et al.\ (2016). ``Mastering the Game of Go with Deep Neural Networks and Tree Search.'' \textit{Nature}, 529, 484--489.
  \item Redelmeier, D.\,H.\ (1981). ``Counting polyominoes: Yet another attack.'' \textit{Discrete Mathematics}, 36(2), 191--203.
\end{enumerate}

\end{document}
